\documentclass{article}

% lualatex handles UTF-8 natively; remove inputenc/fontenc for pdflatex compatibility
\usepackage{hyperref}
\usepackage{amsmath}
\usepackage{mathtools}
\usepackage{amssymb}
\usepackage{amsthm}
\usepackage{dsfont}
\usepackage{stmaryrd}
\usepackage{xcolor}
\usepackage{listings}
\usepackage{microtype}
\usepackage{fullpage}
\usepackage{authblk}
\usepackage{cleveref}

% ---- Theorem environments -----------------------------------------------
\newtheorem{theorem}{Theorem}[section]
\newtheorem{lemma}[theorem]{Lemma}
\newtheorem{definition}[theorem]{Definition}
\newtheorem{example}[theorem]{Example}
\newtheorem{remark}[theorem]{Remark}

% ---- Macros -------------------------------------------------------------
\newcommand{\R}{\mathbb{R}}
\newcommand{\Q}{\mathbb{Q}}
\newcommand{\N}{\mathbb{N}}
\newcommand{\Fin}[1]{\mathsf{Fin}(#1)}
\newcommand{\Flag}[1]{\mathcal{F}^{#1}}
\newcommand{\FlagAlg}[1]{\mathcal{A}^{#1}}
\newcommand{\poshom}[1]{\Phi^{#1}}
\newcommand{\den}[2]{p(#1,\,#2)}
\newcommand{\tdensity}[2]{\pi(#1;\,#2)}
\newcommand{\lean}[1]{\texttt{#1}}
\newcommand{\leanfmt}[1]{\texttt{\small #1}}

% Lean code style
\lstdefinelanguage{Lean4}{
  keywords={def,theorem,lemma,instance,structure,class,import,open,namespace,
             end,where,fun,let,have,show,by,match,with,if,then,else,
             return,do,for,in,noncomputable,abbrev,variable,section,
             native_decide,decide,norm_num,simp,ring,linarith,omega,
             exact,apply,rw,calc,intro,constructor,use,refine,obtain,
             rcases,cases,induction,assumption,contradiction,trivial,rfl},
  keywordstyle=\color{blue!70!black}\bfseries,
  comment=[l]{--},
  commentstyle=\color{gray}\itshape,
  stringstyle=\color{orange!80!black},
  morestring=[b]",
  sensitive=true,
  basicstyle=\ttfamily\small,
  breaklines=true,
  columns=fullflexible,
  escapeinside={(*@}{@*)},
  literate=
    {->}{{$\to$}}2
    {<-}{{$\leftarrow$}}2
    {<=}{{$\leq$}}1
    {>=}{{$\geq$}}1
    {alpha}{{$\alpha$}}1
    {sigma}{{$\sigma$}}1
    {forall}{{$\forall$}}1
    {exists}{{$\exists$}}1
    {phi}{{$\phi$}}1
    {psi}{{$\psi$}}1
}
\lstset{language=Lean4, frame=single, framesep=4pt,
        rulecolor=\color{black!20}, backgroundcolor=\color{gray!5}}

% =========================================================================
\title{Formalizing Flag Algebras in Lean~4 via Computational Reflection}
\date{}

\author[1]{Gyeongwon Jeong}
\affil[1]{School of Computing, KAIST, South Korea}

\author[2]{Seonghun Park}
\affil[2]{School of Computing, KAIST, South Korea}

\author[3]{Jihoon Hyun}
\affil[3]{School of Computing, KAIST, South Korea}

\author[4]{Sang-il Oum}
\affil[4]{Discrete Mathematics Group, Institute for Basic Science (IBS), South Korea}

\author[5]{Hongseok Yang}
\affil[5]{School of Computational Sciences, Korea Institute for Advanced Study (KIAS), South Korea}

% =========================================================================
\begin{document}
\maketitle

% =========================================================================
\begin{abstract}
Razborov's flag algebra method is one of the most powerful tools in extremal
combinatorics, having resolved many open problems about asymptotic subgraph
densities.  Applying it in practice, however, requires combining abstract algebraic
and measure-theoretic reasoning with large external computations: subgraph
density tables computed by enumeration, and semidefinite programming (SDP)
certificates found by numerical solvers.  Formalizing such proofs in a proof
assistant is therefore a challenge on two fronts: the abstract mathematical
structures must be faithfully encoded in a type theory, and the external
computational components must be imported and verified in a way that does not
compromise the overall proof's trustworthiness.

We present a Lean~4 formalization of Razborov's flag algebra method for
graphs, organized around a \emph{reflection-based} architecture.  At the
abstract level, we define flags as quotient types under graph isomorphism,
construct the flag algebra as a quotient module, equip it with a semantic
ordering via positive homomorphisms, and connect flag algebra inequalities to
combinatorial Turán densities through a measure-theoretic framework.  At the
computational level, we introduce a concrete, decidably-equal graph
representation (\lean{Sym2Graph}) and prove adequacy theorems connecting it to
the abstract definitions; this allows \lean{native\_decide} and \lean{decide
+kernel} to discharge hundreds of density and matrix-equality obligations
automatically.  Structural proof obligations---normalizing linear combinations
of flag terms and applying expansion and multiplication identities---are handled
by a suite of custom Lean~4 elaboration tactics that inspect the AST and exploit
a canonical naming convention for flags as a machine-readable encoding of
mathematical structure.

As results, we give complete formal proofs of Mantel's theorem and of the
Erd\H{o}s pentagon theorem ($\pi(K_3; C_5) = 24/625$), the latter including
both the upper bound via a formally verified SDP certificate and the lower
bound via an explicit blow-up construction.  To our knowledge, this is the
first formalization of the flag algebra method in any proof assistant.
\end{abstract}

% =========================================================================
\section{Introduction}
\label{sec:intro}

\paragraph{Flag algebras in extremal combinatorics.}
In extremal graph theory, one frequently asks: among all large graphs on $n$
vertices that avoid some fixed graph $H$ as an induced subgraph, what is the
maximum possible density of another fixed graph $F$?  Such questions are
captured by the \emph{generalized Turán density} $\tdensity{F}{H}$, which is
the limit as $n\to\infty$ of the normalized extremal count.

Razborov's flag algebra method~\cite{razborov2007flag} provides a systematic
framework for deriving upper bounds on these densities.  It works by
constructing a graded, commutative algebra---the \emph{flag algebra}---whose
elements represent linear combinations of induced subgraph patterns, with a
product encoding simultaneous occurrence.  The key insight is that certain
elements of this algebra are guaranteed to be semantically non-negative (their
value under every ``positive homomorphism'' is $\geq 0$), and a certificate of
non-negativity for a carefully chosen element directly yields an upper bound on
the target density.

In practice, those non-negativity certificates are produced by semidefinite
programming: one seeks a positive semidefinite matrix $Q$ such that
$f - \lambda \cdot 1 = \sum_{i,j} Q_{ij} \cdot e_i e_j$ in the flag algebra
(where the $e_i$ are flag basis elements).  Solvers find $Q$ numerically, after
which one must \emph{verify} that $Q$ is indeed positive semidefinite and that
the algebraic identity holds exactly.

\paragraph{The formalization challenge.}
Formalizing a flag algebra proof requires confronting three distinct categories
of proof obligation:

\begin{enumerate}
  \item \textbf{Abstract structure.}  Flags, the flag algebra, positive
    homomorphisms, the semantic cone, the Turán density, and the transfer
    theorems connecting algebra to combinatorics.  These are conceptually
    nontrivial but finite in number.

  \item \textbf{Data-heavy computation.}  For each pair of flags $(F, G)$
    needed in the proof, one must certify the exact rational density
    $\den{F}{G}$.  For a theorem like the Erd\H{o}s pentagon problem, this
    means hundreds of density values over graphs with up to 5 vertices, and
    similarly for flag multiplication tables.  These are computed externally
    and must be imported and verified.

  \item \textbf{Algebraic bookkeeping.}  Flag algebra proofs involve
    manipulating linear combinations of hundreds of flag terms: expanding a
    small flag at a larger size, computing products of typed flags, and
    normalizing sums into a canonical form.  Each step is routine but
    individually tedious.
\end{enumerate}

We address each category with a dedicated technique.  The abstract structure is
encoded directly in Lean~4's dependent type system.  Data-heavy computation is
handled by a \emph{reflection} architecture: we define a decidably-computable
concrete representation of graphs, prove adequacy theorems connecting it to the
abstract definitions, and let \lean{native\_decide}/\lean{decide +kernel}
evaluate the computation inside the proof.  Algebraic bookkeeping is handled by
\emph{custom elaboration tactics} that inspect the proof state's AST,
exploit a canonical naming convention for flag constants, and apply sequences
of domain-specific lemmas automatically.

\paragraph{Contributions.}
\begin{itemize}
  \item \textbf{Abstract formalization.}  We formalize the full abstract
    structure of Razborov's flag algebra in Lean~4: flags as quotients of
    labeled graphs under isomorphism (\S\ref{sec:abstract}), the flag algebra as
    a quotient module with a commutative ring structure
    (\S\ref{sec:abstract}), the semantic ordering via positive homomorphisms, and
    a measure-theoretic framework for the forbidden-subgraph condition
    (\S\ref{sec:forbidden}).

  \item \textbf{Reflection architecture for density and SDP.}  We introduce
    \lean{Sym2Graph}, a finitely-representable graph type, and prove adequacy
    theorems equating abstract flag densities to computable \lean{Sym2Graph}
    densities.  We also verify SDP certificates via LDL$^\top$ decomposition
    checked by \lean{decide +kernel} over $\Q$ (\S\ref{sec:reflection}).

  \item \textbf{Tactic automation.}  We implement a suite of custom Lean~4
    elaboration tactics for linear normalization (\lean{ac\_sort\_pipeline})
    and for systematically applying flag expansion and multiplication identities
    (\lean{prove\_flag\_expand\_with\_forbidden\_flag},
    \lean{prove\_flag\_mul\_with\_forbidden\_flag}) (\S\ref{sec:tactics}).

  \item \textbf{Verified results.}  We give formally complete proofs of
    Mantel's theorem and of $\tdensity{K_3}{C_5} = 24/625$ (the Erd\H{o}s
    pentagon theorem), including both bounds (\S\ref{sec:results}).
\end{itemize}

% =========================================================================
\section{Background: Flag Algebras}
\label{sec:background}

\paragraph{Types and flags.}
Fix a finite simple graph $\sigma$ on vertex set $[k] = \{1,\ldots,k\}$,
called a \emph{type} of size $k$.  A \emph{$\sigma$-flag} is a pair
$(G, \theta)$ where $G$ is a finite simple graph and
$\theta : [k] \hookrightarrow V(G)$ is an injective graph homomorphism from
$\sigma$ to $G$ (i.e., an embedding of the type).  Two $\sigma$-flags
$(G,\theta)$ and $(G',\theta')$ are \emph{isomorphic} if there is a graph
isomorphism $\phi : G \to G'$ with $\phi \circ \theta = \theta'$.  We write
$\Flag{\sigma}_n$ for the set of isomorphism classes of $\sigma$-flags of size
$n$ (necessarily finite), and $\Flag{\sigma} = \bigcup_n \Flag{\sigma}_n$.

The \emph{empty type} $\emptyset$ (size $k=0$) gives ordinary graphs up to
isomorphism.  For this type, flags are just isomorphism classes of graphs.

\paragraph{Flag densities.}
For $\sigma$-flags $F$ of size $m$ and $G$ of size $n \geq m$, the
\emph{induced density} $\den{F}{G}$ is the probability that a uniformly random
$(m - k)$-subset of $V(G) \setminus \operatorname{Im}(\theta_G)$ induces a
copy of $F$ compatible with the type embedding.  More precisely:
\[
  \den{F}{G} \;=\;
  \frac{\bigl|\{\text{injections } \iota : V(F) \to V(G) \mid \iota \text{ preserves type and induces } F\}\bigr|}
       {\binom{n-k}{m-k} \cdot (m-k)!}.
\]
This quantity is always in $[0,1]$ and is a rational number whenever $G$ is
finite.

\paragraph{The flag algebra.}
Fix a type $\sigma$.  The \emph{flag algebra} $\FlagAlg{\sigma}$ is
constructed as follows.  Let $\mathbb{R}[\Flag{\sigma}]$ be the free
$\R$-module with basis $\Flag{\sigma}$.  Define the \emph{zero space}
$\mathcal{Z}^\sigma$ as the subspace generated by elements of the form
\[
  F - \sum_{G \in \Flag{\sigma}_n} \den{F}{G} \cdot G,
\]
for each $F \in \Flag{\sigma}_m$ and each $n \geq m$.  Then
$\FlagAlg{\sigma} = \mathbb{R}[\Flag{\sigma}] / \mathcal{Z}^\sigma$.

The product of $F \in \Flag{\sigma}_m$ and $F' \in \Flag{\sigma}_{m'}$ at size
$\ell \geq m + m' - k$ is
\[
  F \cdot F' = \sum_{G \in \Flag{\sigma}_\ell} \den{F,F'}{G} \cdot G,
\]
where $\den{F,F'}{G}$ is the probability that two independently chosen
(compatible) subsets of $V(G)$ induce copies of $F$ and $F'$ respectively.
This product is well-defined on the quotient and makes $\FlagAlg{\sigma}$ into
a commutative $\R$-algebra.

\paragraph{The semantic cone and positivity.}
A \emph{positive homomorphism} is an $\R$-algebra homomorphism
$\phi : \FlagAlg{\emptyset} \to \R$ satisfying $\phi(F) \geq 0$ for all
flags $F$.  It can be shown that positive homomorphisms are in bijection with
convergent sequences of graphs: every infinite graph sequence that
``converges'' (in the sense that all flag densities have limits) determines a
unique positive homomorphism.

The \emph{semantic cone} is $\{f \in \FlagAlg{\sigma} \mid \forall \phi, \phi(f) \geq 0\}$.
Razborov's key theorem is:
\begin{theorem}[Razborov~\cite{razborov2007flag}]
  If $A$ is a positive semidefinite matrix and $e_1,\ldots,e_r$ are typed flags
  in $\FlagAlg{\sigma}$, then $\sum_{i,j} A_{ij} \cdot e_i e_j$ lies in the
  semantic cone after averaging over the type embedding (the ``downward'' operator
  $\llbracket \cdot \rrbracket_\sigma$).
\end{theorem}
This is the main mechanism by which SDP certificates produce valid
non-negativity proofs.

\paragraph{Turán density and the pentagon problem.}
For unlabeled graphs $F$ and $H$, the \emph{generalized Turán density}
$\tdensity{F}{H}$ is:
\[
  \tdensity{F}{H} = \lim_{n\to\infty} \frac{1}{\binom{n}{|V(F)|}}
  \max\bigl\{|\text{copies of } F \text{ in } G \mid G \text{ is } H\text{-free},\; |V(G)|=n\bigr\}.
\]
The \emph{Erd\H{o}s pentagon problem} asks for $\tdensity{K_3}{C_5}$: the
maximum density of triangles in triangle-free-with-respect-to-$C_5$-free
graphs.  The answer, $24/625$, was proved by Grzesik~\cite{grzesik2012} and
independently by Hatami, Hladk\'y, Kr\'al, Norine, and
Razborov~\cite{hatami2012}; it is achieved by the blow-up of $C_5$.

% =========================================================================
\section{Abstract Formalization}
\label{sec:abstract}

\subsection{Flags as Quotient Types}

In Lean~4 we exploit dependent types and quotient types to encode the flag
algebra faithfully.  The central definitions live in
\lean{LeanFlagAlgebras.FlagAlgebra.FlagDef}.

A \emph{type graph} of size $k$ is a \lean{SimpleGraph (Fin k)}.  A
\emph{labeled graph} over type $\sigma$ with vertex set $V$ is a pair
\lean{(graph : SimpleGraph V, type\_embed : $\sigma$ $\hookrightarrow_g$ graph)},
where the embedding records a copy of the type inside the host graph.  Two
labeled graphs over the same type are \emph{flag-isomorphic} (\lean{$\sim$f})
if they are connected by a graph isomorphism that preserves the type embedding.
A \emph{flag} is then the quotient:

\begin{lstlisting}
def Flag (σ : FlagType (Fin n₀)) (V : Type) : Type :=
  Quotient (labeledGraphSetoid σ V)
\end{lstlisting}

We write \lean{FlagWithSize $\sigma$ $n$} for \lean{Flag $\sigma$ (Fin $n$)},
and \lean{FinFlag $\sigma$ := $\Sigma$ ($n$ : $\N$), FlagWithSize $\sigma$ $n$}
for the disjoint union over all sizes.  These types are automatically finite
(for each fixed size) and countable (in total), inheriting these properties from
Lean's standard library.

\subsection{The Flag Algebra as a Quotient Module}

The flag algebra is built in \lean{LeanFlagAlgebras.FlagAlgebra.FlagAlgebra}.
We form the free $\R$-module of finitely-supported functions:
\begin{lstlisting}
abbrev FlagVector σ := FinFlag σ →₀ ℝ
\end{lstlisting}
and define the \emph{zero space} \lean{ZeroSpace $\sigma$} as the submodule
generated by the density-expansion relations.  The flag algebra is:
\begin{lstlisting}
def FlagAlgebra σ := FlagVector σ ⧸ ZeroSpace σ
\end{lstlisting}
We prove that this carries a commutative $\R$-algebra structure, where the
product $[F] \cdot [F']$ is defined via the flag multiplication
\lean{flagMulWithSize}.

\subsection{Semantic Ordering via Positive Homomorphisms}

The ordering on \lean{FlagAlgebra $\sigma$} is semantic: we say $f \leq g$ if
$g - f$ lies in the \emph{semantic cone}, defined as
\begin{lstlisting}
def semanticCone σ : Set (FlagAlgebra σ) :=
  {f | ∀ φ : PositiveHom σ, φ f ≥ 0}
\end{lstlisting}
where \lean{PositiveHom $\sigma$} is the type of $\R$-algebra homomorphisms
$\phi : \lean{FlagAlgebra}\ \sigma \to \R$ that map every unit vector
(flag basis element) to a non-negative real.  Proving that this ordering is a
preorder compatible with the algebra operations is straightforward from the
definitions.

The positive semidefinite quadratic form result is formalized in
\lean{LeanFlagAlgebras.FlagAlgebra.QuadraticForm}:
\begin{lstlisting}
theorem flagQuadraticForm_nonneg
    (M : Matrix (Fin r) (Fin r) ℝ) (hM : M.PosSemidef)
    (v : Fin r → FlagAlgebra σ) :
    0 ≤ flagQuadraticForm M v
\end{lstlisting}
This is the key lemma that converts a PSD certificate into a membership in the
semantic cone.

\subsection{The Forbidden-Subgraph Framework}
\label{sec:forbidden}

A novel aspect of our formalization is a \emph{general forbidden-subgraph
reasoning rule} that works inside the ambient theory of simple graphs rather
than requiring a problem-specific axiomatization, as in Razborov's original
presentation.

The space of positive homomorphisms $\poshom{\sigma}$ is given a topology (as a
subspace of $\R^{\lean{FinFlag}\ \sigma}$, which is compact by Tychonoff), and
each positive homomorphism $\phi_0 : \lean{FlagAlgebra}\ \emptyset \to \R$
induces a probability measure $\mathbb{P}^{\phi_0}$ on $\poshom{\sigma}$ (the
``random flag extension'' measure).  This is formalized in
\lean{LeanFlagAlgebras.FlagAlgebra.FlagSequence} and
\lean{LeanFlagAlgebras.FlagAlgebra.RandomHom}.

We then define:
\begin{lstlisting}
def forbidLE (F_forbid : FinFlag ∅ₜ) (f g : FlagAlgebra σ) : Prop :=
  ∀ (φ₀ : PositiveHom ∅ₜ), φ₀ ⟦F_forbid⟧ = 0 →
    ℙ[φ₀] {φ | φ f ≤ φ g} = 1
\end{lstlisting}
The notation $f \leq_{[\mathtt{C5}]} g$ means: almost surely (under the
measure induced by any positive homomorphism that gives $C_5$ density zero),
$\phi(f) \leq \phi(g)$.  The key transfer theorem is:
\begin{lstlisting}
theorem generalizedTuranDensity_le_of_forbidLE
    (h : 0 < c) (hineq : F_target ≤[H_forbid] c • 1) :
    generalizedTuranDensity F_target H_forbid ≤ c
\end{lstlisting}
This theorem, proved in \lean{LeanFlagAlgebras.Forbid.Basic}, is the bridge
between the flag algebra world and the combinatorial Turán density.  It is
completely general: any flag algebra inequality proved under the
forbidden-subgraph assumption yields a Turán density bound for any forbidden
graph $H$.

% =========================================================================
\section{The Reflection Layer}
\label{sec:reflection}

The abstract formalization in \S\ref{sec:abstract} gives the right mathematical
definitions but is not directly evaluable: \lean{flagDensity} is defined via
non-constructive quotients, and \lean{Decidable} instances are not available
directly.  The reflection layer bridges this gap.

\subsection{Proof by Reflection in ITP}

The proof-by-reflection pattern in interactive theorem proving has three
components: (1) a \emph{concrete representation} over which a decision
procedure can evaluate; (2) an \emph{adequacy theorem} proving that the
decision procedure is sound with respect to the abstract definition; and (3)
a \emph{decision procedure} (the evaluator) that, given the concrete representation,
produces a Boolean answer which \lean{decide} converts into a proof.  Our
architecture applies this pattern at multiple granularities.

\subsection{Concrete Graph Representations}

We introduce \lean{Sym2Graph n} in
\lean{LeanFlagAlgebras.FlagAlgebra.Compute.Basic}:
\begin{lstlisting}
structure Sym2Graph (n : ℕ) where
  edges : Finset (Sym2 (Fin n))
  edges_valid : ∀ e ∈ edges, ¬e.IsDiag
\end{lstlisting}
This represents a graph on $\{0,\ldots,n-1\}$ as a finite set of unordered
non-diagonal pairs.  Unlike \lean{SimpleGraph (Fin n)} (which uses a
\lean{Prop}-valued adjacency relation), every operation on \lean{Sym2Graph n}
is computable: membership in \lean{edges} is decidable by \lean{Finset}
membership.

The same file derives the \lean{Fintype} instances needed to make graph
isomorphism decidable:
\begin{lstlisting}
instance : Fintype (G₁ ↪g G₂) := ...   -- graph embeddings
instance : Fintype (G₁ ≃g G₂) := ...   -- graph isomorphisms
instance : Decidable (Nonempty (G ≃f G')) := ...  -- flag isomorphism
\end{lstlisting}
These instances are engineered carefully: they enumerate candidates (all
injections from $V(G_1)$ to $V(G_2)$, or all bijections) and filter those
that preserve adjacency and type embeddings.  This reduces the question ``are
these two labeled graphs isomorphic?'' to a finite, decidable search.

\subsection{Adequacy Theorems}

The connection between the abstract and concrete worlds is established by
adequacy theorems in \lean{LeanFlagAlgebras.FlagAlgebra.Compute.FlagDensity}:

\begin{lstlisting}
theorem flagDensity₁_eq_sym2EmptyTypeFlagDensity₁
    (F : Sym2EmptyTypedFlag m) (G : Sym2EmptyTypedFlag n) :
    flagDensity₁ F.toFlag G.toFlag =
    sym2EmptyTypeFlagDensity₁ F G

theorem flagDensity₂_eq_sym2EmptyTypeFlagDensity₂
    (F₀ : Sym2EmptyTypedFlag m₀) (F₁ : Sym2EmptyTypedFlag m₁)
    (G : Sym2EmptyTypedFlag n) :
    flagDensity₂ F₀.toFlag F₁.toFlag G.toFlag =
    sym2EmptyTypeFlagDensity₂ F₀ F₁ G
\end{lstlisting}
The right-hand sides---\lean{sym2EmptyTypeFlagDensity₁} and
\lean{sym2EmptyTypeFlagDensity₂}---are defined entirely over \lean{Sym2Graph}
and are fully computable.  After rewriting a density goal with these theorems,
the goal becomes a decidable proposition about a finite computation over
\lean{Finset}, which \lean{native\_decide} can evaluate directly.

\subsection{Elaboration-Time Theorem Generation}

Having established the adequacy theorems, we need to generate and prove the
density and multiplication lemmas needed for each specific theorem.  Rather
than writing these by hand, we use Lean~4's elaboration monad to generate them
automatically from precomputed data.

The elaboration macro \lean{load\_flag\_pair\_density\_theorems} in
\lean{LeanFlagAlgebras.ErdosPentagon.Densities.DensityLoader} reads a JSON
file containing density values, and for each entry \lean{(F1, F2, G, p/q)}
generates and immediately proves the theorem:
\begin{lstlisting}
@[simp]
theorem flagDensity₂_Flag_5_1_0_3_Flag_5_1_0_7_Flag_5_0_0_2
    : flagDensity₂ Flag_5_1_0_3 Flag_5_1_0_7 Flag_5_0_0_2 = 3/10
  := by
  delta Flag_5_1_0_3 Flag_5_1_0_7 Flag_5_0_0_2
  rw [flagDensity₂_eq_sym2FlagDensity₂]
  native_decide
\end{lstlisting}
The three proof steps are: unfold the concrete flag definitions (making their
\lean{Sym2Graph} structure explicit), rewrite with the adequacy theorem, then
call \lean{native\_decide} to evaluate the resulting finite computation.
Hundreds of such theorems are generated and proved at compile time.  The same
elaboration pattern is used by \lean{MulLoader} for flag multiplication tables.

\subsection{SDP Certificate Verification via LDL$^\top$}
\label{sec:sdp}

The upper bound proof for the Erd\H{o}s pentagon theorem requires showing
that a sum of three quadratic forms (one for each 1-vertex type) is
non-negative under the $C_5$-free assumption.  The certificates are three
positive semidefinite matrices $P$, $Q$, $R$ over $\Q$ (each $8 \times 8$),
found externally by an SDP solver.

We verify positive semidefiniteness via an LDL$^\top$ decomposition.  For
matrix $P$, we exhibit a lower-triangular matrix $L_P$ and a diagonal matrix
$D_P$ with non-negative diagonal entries $d_P$, defined as explicit rational
constants in \lean{LeanFlagAlgebras.ErdosPentagon.Matrix.PosSemiDef}:
\begin{lstlisting}
def P : Matrix (Fin 8) (Fin 8) ℚ := !![(24/625 : ℚ), (-36/625 : ℚ), ...]
def dP : Fin 8 → ℚ := ![(24/625 : ℚ), (223/625 : ℚ), ..., 0, 0, 0, 0]
def LP : Matrix (Fin 8) (Fin 8) ℚ := !![(1 : ℚ), 0, ...; (-3/2 : ℚ), 1, ...]
\end{lstlisting}
Positive semidefiniteness then follows from two lemmas:
\begin{lstlisting}
lemma dP_nonneg (i : Fin 8) : 0 ≤ dP i := by
  fin_cases i <;> norm_num [dP]

lemma P_eq_LDL : P = LP * Matrix.diagonal dP * LPᵀ := by
  decide +kernel

theorem P_posSemidef : P.PosSemidef :=
  posSemidef_of_eq_mul_diagonal_mul_transpose dP_nonneg P_eq_LDL
\end{lstlisting}
The matrix equality \lean{P\_eq\_LDL} is a decidable proposition over $\Q$
(rational matrix equality), verified by \lean{decide +kernel}---the Lean
kernel evaluates the matrix product symbolically and checks equality entry by
entry.  This uses no compiled native code, so no axiom beyond the standard
Lean kernel is assumed.

\paragraph{The trust hierarchy.}
We deliberately use two different decision procedures with different trust
profiles:
\begin{itemize}
  \item \lean{native\_decide} compiles the goal to native code and evaluates
    it.  It relies on the correctness of the Lean-to-native compiler, which is
    outside the kernel.  We use it for density tables, where the computation is
    large but the individual claims are simple (a rational equality between the
    defined density and a stated value).

  \item \lean{decide +kernel} evaluates inside the kernel.  It is slower but
    uses no trusted computation outside the Lean kernel itself.  We use it for
    the SDP matrix equalities, where trust is most critical.
\end{itemize}
In both cases, the external computation (the density enumeration scripts, the
SDP solver) is responsible only for producing \emph{candidates}---values to
claim.  The Lean proof is responsible for verifying each claim against the
formal definition.

% =========================================================================
\section{The Tactic Layer}
\label{sec:tactics}

The reflection layer handles data-heavy computation.  A complementary layer of
custom elaboration tactics handles the structurally repetitive proof steps that
arise in every flag algebra argument.

\subsection{Naming Convention as Machine-Readable Encoding}

A key enabling design decision is that every flag and flag algebra constant is
named according to the scheme \lean{Flag\_n\_k\_m\_i} and
\lean{FlagAlgebra\_n\_k\_m\_i}, where $n$ is the vertex count, $k$ and $m$
identify the type graph, and $i$ is the canonical index within that
isomorphism class.  This convention is not merely organizational: it serves as
a machine-readable encoding of mathematical structure that all custom tactics
exploit.  By parsing the trailing integers from a constant's name, a tactic can
determine:
\begin{itemize}
  \item the \emph{canonical order} of a flag term in a linear combination,
  \item the \emph{precomputed lemma name} for its density or multiplication
    identity, and
  \item the \emph{flag type parameters} needed to instantiate the appropriate
    adequacy theorem.
\end{itemize}
This design is similar to the use of canonical forms in certified compilation:
the naming scheme is a ``normal form'' that tactics can compute with directly,
without needing to reduce the mathematical content of the expression.

\subsection{Linear Normalization: \lean{ac\_sort\_pipeline}}

\paragraph{The problem.}
After expanding a flag algebra element at a larger size, the proof obligation
is typically an equality between two linear combinations of flag terms that are
provably equal but syntactically in different orders.  Generic tactics such as
\lean{ring} or \lean{ac\_rfl} handle this in principle but become prohibitively
slow on large expressions ($\geq 20$ terms) because they must search over all
permutations of summands.

\paragraph{The solution.}
The \lean{sort} family of tactics, implemented in
\lean{LeanFlagAlgebras.ErdosPentagon.SortTactic}, normalizes a linear
combination by:
\begin{enumerate}
  \item \emph{Flattening.}  The tactic traverses the expression's AST,
    pattern-matching on \lean{HAdd.hAdd}, \lean{HSub.hSub},
    \lean{HSMul.hSMul}, and \lean{Neg.neg}, to produce a flat list of
    \lean{(base, coefficient)} pairs.
  \item \emph{Sorting.}  Each base term is assigned a sort key: the trailing
    integer in its constant name (e.g., \lean{FlagAlgebra\_5\_1\_0\_3} gets
    key $3$), with the pretty-printed string as a tiebreaker.  The list is
    sorted by this key using insertion sort.
  \item \emph{Rebuilding.}  The sorted list is reassembled into a Lean
    expression (left-associative sum of scalar multiples).
  \item \emph{Equality proof.}  The tactic proves that the original and sorted
    expressions are equal by \lean{ac\_rfl} or \lean{abel\_nf}, then replaces
    the goal with the sorted form.
\end{enumerate}
The full pipeline \lean{ac\_sort\_pipeline} additionally runs \lean{norm\_num}
before and after sorting, and collects like terms with \lean{$\leftarrow$ add\_smul}
after sorting, so that the resulting expression is fully simplified.

\paragraph{Why this is necessary, not cosmetic.}
The comment in \lean{ErdosPentagon.Lemmas} is revealing:
``\texttt{-- sort\_at -- takes more than 10 minutes}''.
The generic \lean{sort\_at} (which uses \lean{simp} without index guidance)
times out on expressions with $\sim 25$ flag terms.  The index-guided
\lean{ac\_sort} completes the same step in under a second, by exploiting domain
knowledge to bypass combinatorial search.

\subsection{Flag Expansion and Multiplication Tactics}

\paragraph{The problem.}
A central step in every flag algebra proof is the \emph{expansion} identity: a
small flag $F \in \Flag{\sigma}_m$, viewed in the algebra at size $n > m$,
equals a weighted sum of all $n$-vertex flags.  Similarly, the product of two
typed flags equals a linear combination at the product size.  These identities
must be verified for each specific flag and expansion size that appears in the
proof---dozens of instances per theorem.

\paragraph{\lean{prove\_flag\_expand\_with\_forbidden\_flag} N.}
This tactic, implemented in \lean{LeanFlagAlgebras.Logic.Tactic}, proves goals
of the form:
\[
  \bigl(\lean{FlagAlgebra\_3\_0\_0\_3} =_{a} 0\bigr)
  \;\vdash_a\;
  \lean{FlagAlgebra\_2\_0\_0\_1} =_{a} \tfrac{1}{3} \cdot \lean{FlagAlgebra\_3\_0\_0\_1}
    + \tfrac{2}{3} \cdot \lean{FlagAlgebra\_3\_0\_0\_2}
\]
where $=_{a}$ is the assertion operator in the logical framework.  The tactic:
\begin{enumerate}
  \item Scans the goal's AST for all \lean{FlagAlgebra\_*} and \lean{Flag\_*}
    constants using \lean{collectPrefixConstants}, and parses their indices
    $(n, k, m, i)$ using \lean{parseFlagAlgebraIndices?}.
  \item Unfolds the flag definitions by name, introduces the positive
    homomorphism $\phi$ and the forbidden-flag hypothesis $h$.
  \item Applies \lean{unitVector\_quot\_eq\_sum} to expand the flag at size
    $N$, then rewrites with the preloaded lemmas \lean{flagSet\_N\_k\_m\_eq\_univ}
    and \lean{flagSet\_N\_k\_m\_val\_eq} (themselves generated by the loaders).
  \item Simplifies the expansion using $h$ (the forbidden flag has value~0)
    and closes with \lean{ring\_nf}.
\end{enumerate}

\paragraph{\lean{prove\_flag\_mul\_with\_forbidden\_flag} N.}
This tactic handles product identities under a forbidden-flag assumption.  It
follows the same structure but applies \lean{unitVector\_quot\_mul\_eq\_flagMul\_quot}
and unfolds the flag multiplication at size $N$.  It also normalizes the
multiplication order: if the two operands' indices are in non-canonical order,
it first applies \lean{mul\_comm} before the rest of the proof, ensuring that
precomputed lemmas (which assume a fixed canonical ordering) can be applied
directly.

The result of these two tactics is that expansion and multiplication identities
that would individually require 15--20 tactic steps are each proved by a single
tactic invocation:
\begin{lstlisting}
example : FlagAlgebra_3_0_0_3 =ₐ 0
    ⊢ₐ FlagAlgebra_2_0_0_1 =ₐ (1/3 : ℝ) • FlagAlgebra_3_0_0_1
                             + (2/3 : ℝ) • FlagAlgebra_3_0_0_2
  := by prove_flag_expand_with_forbidden_flag 3
\end{lstlisting}

% =========================================================================
\section{Results}
\label{sec:results}

\subsection{Mantel's Theorem}

Mantel's theorem states that in a triangle-free graph on $n$ vertices,
the number of edges is at most $\lfloor n^2/4 \rfloor$ (achieved by
the complete bipartite graph $K_{\lfloor n/2\rfloor, \lceil n/2\rceil}$).
In flag algebra terms, the density of edges is at most $1/2$ whenever the
density of triangles is $0$:

\begin{lstlisting}
theorem Mantel_theorem
    : K2 ≤ (1/2 : ℝ) • 1 + K3
\end{lstlisting}

Here $K_2$ and $K_3$ denote the elements of \lean{FlagAlgebra $\emptyset_t$}
corresponding to the single edge and the triangle, respectively.  The proof uses
an explicit square in the algebra of 1-vertex-typed flags:
\[
  \bigl\llbracket(\lean{FlagAlgebra\_2\_1\_0\_0} - \lean{FlagAlgebra\_2\_1\_0\_1})^2
  \bigr\rrbracket \;\geq\; 0,
\]
which gives the inequality after expansion and simplification.  The SDP
certificate here is small enough that it can be written explicitly (no external
solver needed), making Mantel's theorem a clean illustration of the method
before the heavier machinery is needed for the pentagon problem.

\subsection{The Erd\H{o}s Pentagon Theorem}

The main result is:
\begin{lstlisting}
theorem ErdosPentagon_Turan
    : generalizedTuranDensity K3 C5 = 24/625
\end{lstlisting}
proved as the conjunction of two bounds.

\paragraph{Upper bound.}
The upper bound \lean{ErdosPentagon\_Turan\_upperBound} is proved in one line:
\begin{lstlisting}
theorem ErdosPentagon_Turan_upperBound
    : generalizedTuranDensity K3 C5 ≤ 24/625 :=
  generalizedTuranDensity_le_of_forbidLE (by norm_num)
    ErdosPentagon_flagAlgebra
\end{lstlisting}
The key lemma \lean{ErdosPentagon\_flagAlgebra} establishes the flag algebra
inequality $C_5\text{-density} \leq_{[K_3]} \frac{24}{625} \cdot 1$,
proved by exhibiting three quadratic forms (one per 1-vertex type, with PSD
matrices $P$, $Q$, $R$) that are non-negative under the $K_3$-free assumption,
and then verifying that their sum dominates the $C_5$ density element.  The
PSD check for each matrix uses the LDL$^\top$ approach from \S\ref{sec:sdp}.

\paragraph{Lower bound.}
The lower bound \lean{ErdosPentagon\_Turan\_lowerBound} is proved via an
explicit blow-up construction:
\begin{lstlisting}
def blowUp (G : SimpleGraph V) (n : ℕ) : SimpleGraph (V × Fin n) :=
  { Adj v w := G.Adj v.1 w.1, ... }
\end{lstlisting}
The blow-up of $C_5$ by factor $n$ gives a $K_3$-free graph
(\lean{blowUp\_K3\_free}) on $5n$ vertices containing at least $n^5$ induced
copies of $C_5$ (\lean{subgraphCount\_blowUp\_C5\_ge}).  Dividing by
$\binom{5n}{5}$ and taking $n\to\infty$ gives the lower bound:
\begin{lstlisting}
theorem ErdosPentagon_Turan_lowerBound
    : generalizedTuranDensity K3 C5 ≥ 24/625
\end{lstlisting}
The limit argument uses Lean's \lean{Filter.Tendsto} framework and the fact
that $n^5 / \binom{5n}{5} \to 24/625$ by an exact algebraic computation
(\lean{norm\_num}).

\paragraph{No axioms beyond the kernel.}
The main theorem \lean{ErdosPentagon\_Turan} depends on \lean{native\_decide}
for the density table lemmas (which introduces a dependency on the native-code
compiler) and on \lean{decide +kernel} for the SDP matrix equalities (which
uses only the kernel).  There are no \lean{sorry}s in the proof path of either
result.

% =========================================================================
\section{Related Work}
\label{sec:related}

\paragraph{Formalizations of combinatorics.}
Several combinatorial theorems have been formalized in Lean and related
systems.  Dillies and Mehta~\cite{dillies2022szemeredi} formalized
Szemerédi's Regularity Lemma in Lean~4; Mehta~\cite{mehta2022kruskal}
formalized Kruskal-Katona; Subercaseaux et al.~\cite{subercaseaux2024hexagon}
formally verified the empty hexagon number using a SAT-based pipeline.  Our
work differs in that the central method (flag algebras) relies on an algebraic
quotient construction and measure-theoretic limits, rather than purely
combinatorial enumeration.

\paragraph{Computer-assisted flag algebra arguments.}
The Flagmatic software system~\cite{razborov2010flagmatic} automates the
computation of flag algebra certificates, but produces results that are
validated only informally.  Our work provides the first formally verified
end-to-end pipeline from flag algebra certificates to combinatorial theorems.

\paragraph{Proof by reflection.}
Proof by reflection is a classical technique in the Coq community (see
Chlipala~\cite{chlipala2013cpdt} for a survey), and has been used in
Lean for verified computation (e.g., in Mathlib's \lean{decide} and
\lean{norm\_num} tactics).  Our use of \lean{Sym2Graph} as a concrete mirror
of \lean{SimpleGraph} follows the same pattern, extended to the domain of flag
density computation.

\paragraph{Tactic meta-programming.}
Custom tactics for domain-specific proof automation have been developed in many
contexts~\cite{ebner2017structured,sozeau2008coq}.  Our tactics exploit the
naming convention as a form of \emph{reflected metadata}: the name of a
constant encodes information that the tactic exploits without having to reduce
the mathematical content of the expression.  This is related to, but distinct
from, approaches based on type-class inference or user-facing attribute
annotations.

% =========================================================================
\section{Conclusion}
\label{sec:conclusion}

We have presented the first formalization of Razborov's flag algebra method in
a proof assistant, realized in Lean~4.  The formalization is organized around
a two-layer architecture.  The \emph{reflection layer} connects the abstract
algebraic definitions to a concrete, decidably-computable graph representation
via adequacy theorems, enabling \lean{native\_decide} and \lean{decide +kernel}
to automatically discharge hundreds of density and SDP certificate obligations.
The \emph{tactic layer} provides custom elaboration tactics---exploiting a
canonical naming scheme for flag constants---that handle the structural
bookkeeping (linear normalization, expansion identities, multiplication
identities) of flag algebra proofs.

The end results are formally complete proofs of Mantel's theorem and the
Erd\H{o}s pentagon theorem ($\tdensity{K_3}{C_5} = 24/625$).  We believe
the architecture generalizes: the $\sim 100$ known flag algebra results in
extremal combinatorics all involve the same three categories of proof
obligation (abstract structure, data-heavy computation, and algebraic
bookkeeping), and our framework provides reusable infrastructure for all three.

\paragraph{Future work.}
Replacing \lean{native\_decide} with \lean{decide +kernel} throughout (or
with a formally verified external checker) would eliminate the remaining
dependency on the native compiler.  Extending the framework beyond graphs---to
hypergraphs, directed graphs, or other combinatorial structures---would require
generalizing the type-parameter conventions but no new conceptual machinery.
Automating the discovery of SDP certificates from within Lean (rather than
importing them from external solvers) remains a longer-term goal.

% =========================================================================
\bibliographystyle{plain}
\begin{thebibliography}{99}

\bibitem{razborov2007flag}
A.~A. Razborov.
\newblock Flag algebras.
\newblock \textit{Journal of Symbolic Logic}, 72(4):1239--1282, 2007.

\bibitem{grzesik2012}
A.~Grzesik.
\newblock On the maximum number of five-cycles in a triangle-free graph.
\newblock \textit{Journal of Combinatorial Theory, Series B},
  102(5):1061--1066, 2012.

\bibitem{hatami2012}
H.~Hatami, J.~Hladk\'{y}, D.~Kr\'{a}l, S.~Norine, and A.~Razborov.
\newblock On the number of pentagons in triangle-free graphs.
\newblock \textit{Journal of Combinatorial Theory, Series A},
  120(3):722--732, 2013.

\bibitem{dillies2022szemeredi}
Y.~Dillies and B.~Mehta.
\newblock Formalising Szemer\'{e}di's regularity lemma in Lean.
\newblock In \textit{Proc.\ ITP 2022}, LIPIcs~237, 2022.

\bibitem{mehta2022kruskal}
B.~Mehta.
\newblock Formalising the Kruskal-Katona theorem in Lean.
\newblock In \textit{Proc.\ ITP 2022}, 2022.

\bibitem{subercaseaux2024hexagon}
B.~Subercaseaux, M.~J.~H. Heule, J.~Mackey, J.~Meadows, R.~Tao, and
  C.~Wu.
\newblock Formal verification of the empty hexagon number.
\newblock In \textit{Proc.\ ITP 2024}, LIPIcs~309, 2024.

\bibitem{razborov2010flagmatic}
A.~A. Razborov.
\newblock On 3-hypergraphs with forbidden 4-vertex configurations.
\newblock \textit{SIAM Journal on Discrete Mathematics}, 24(3):946--963, 2010.

\bibitem{chlipala2013cpdt}
A.~Chlipala.
\newblock \textit{Certified Programming with Dependent Types}.
\newblock MIT Press, 2013.

\bibitem{ebner2017structured}
G.~Ebner, S.~Ullrich, J.~Roesch, J.~Avigad, and L.~de Moura.
\newblock A metaprogramming framework for formal verification.
\newblock \textit{Proc.\ ACM Program.\ Lang.}, 1(ICFP):34:1--34:29, 2017.

\bibitem{sozeau2008coq}
M.~Sozeau and N.~Oury.
\newblock First-class type classes.
\newblock In \textit{Proc.\ TPHOLs 2008}, LNCS~5170, 2008.

\end{thebibliography}

\end{document}
